% !TEX program = lualatex
\documentclass[12pt, a4paper]{article}
\usepackage{master}

\DeclareWorkGR{Cat}{범주론}{Κατηγορίαι}{Categoriae}
\DeclareWorkGR{Top}{토피카}{Τοπικά}{Topica}
\DeclareWorkGR{Met}{형이상학}{Μεταφυσικά}{Metaphysica}
\DeclareWorkGR{APo}{분석론 후서}{Ἀναλυτικὰ ὕστερα}{Analytica posteriora}

\fancyhead[L]{\footnotesize\color{midgray}오르가논 강독}
\fancyhead[R]{\footnotesize\color{midgray}특별 세션 α}

\begin{document}
\thispagestyle{firstpage}

\weeksection{특별 세션 α}{범주론·이사고게 종합 I}{하나의 체계, 열린 문제들, 고대 주석가들}

\subsection{두 텍스트가 함께 이루는 것}

\subsubsection{존재론 위의 술어화 이론}

지난 7주 동안 우리는 두 텍스트를 읽었습니다 --- 아리스토텔레스의 범주론(2--5주차)과 포르피리오스의 이사고게(6--8주차). 이 두 텍스트는 서로 다른 물음에 답합니다. 범주론은 `무엇이 존재하는가'를 묻고(존재론), 이사고게는 `술어가 주어에 어떻게 관계하는가'를 묻습니다(술어화 이론). 그러나 이 두 텍스트는 고대부터 하나의 교육 단위로 읽혔습니다 --- 이사고게를 먼저 읽고, 그 다음에 범주론을 읽는 것이 5세기 이래 신플라톤주의 학교의 표준 교육 과정이었습니다.

에반겔리우(Evangeliou 1996)가 주장하듯이, 이사고게는 범주론의 존재론적 틀이 전제하지만 명시적으로 서술하지 않는 술어화 이론을 제공합니다.\footnote{Christos Evangeliou, \gri{Aristotle's Categories and Porphyry} (Leiden: Brill, 1996).} 범주론은 존재자를 `주어에 대해 말해지는 것'(본질적 술어화)과 `주어 안에 있는 것'(부수적 내재)의 두 축으로 분류합니다. 이사고게는 이 두 축 위에서 작동하는 술어화의 다섯 가지 방식을 체계화합니다.

{\footnotesize
\begin{concepttable}{|l|l|}
\hline
\rowcolor{tableheadblue}
{\color{h1navy}\bfseries 범주론의 사중 분류} &
{\color{h1navy}\bfseries 이사고게의 대응} \\ \hline
(a) 말해짐 + 안에 없음 (제2실체) & 류, 종 \\ \hline
(b) 말해지지 않음 + 안에 없음 (제1실체) & 개별자 (술어가능자 아님) \\ \hline
(c) 말해짐 + 안에 있음 (보편적 속성) & 고유성, 보편적 부수성 \\ \hline
(d) 말해지지 않음 + 안에 있음 (개별적 속성) & 개별적 부수성 \\ \hline
종차: 본질적이지만 실체가 아님 & 사중 분류의 빈자리를 채움 \\ \hline
\end{concepttable}
}

\subsubsection{종차의 문제 --- 사중 분류가 예비하지 못한 자리}

이 대응에서 가장 흥미로운 지점은 종차(\gr{διαφορά})입니다. 7주차에서 살펴보았듯이, 종차는 본질적이면서도(류 + 종차 = 정의) 실체 범주에 속하지 않습니다. 범주론의 사중 분류에서 종차는 (a)에 속할까요? 그렇다면 종차는 제2실체의 일종이 됩니다 --- 그러나 `이성적인'은 종이나 류가 아닙니다. (c)에 속할까요? 그렇다면 종차는 보편적 속성이 됩니다 --- 그러나 종차는 본질적이지 부수적이지 않습니다. 이 딜레마는 범주론이 이사고게 없이는 완전하지 않다는 것을 보여줍니다 --- 포르피리오스가 종차를 별도의 술어가능자로 추가한 것은 범주론 자체의 구조적 빈자리를 메우는 작업이었습니다.

\subsubsection{포르피리오스의 나무와 10범주의 내부 구조}

포르피리오스의 나무(실체 → 물체 → 생명 있는 물체 → 동물 → 이성적 동물 → 사람 → 소크라테스)는 실체 범주의 내부 구조를 보여줍니다. 그러나 원리적으로 다른 범주에도 같은 류-종-종차 위계를 적용할 수 있습니다. 양 범주 안에서 `이산적 양'과 `연속적 양'은 종차에 의한 분할이고(3주차), 질 범주 안에서 `상태', `역량', `감각적 질', `모양'은 네 종류로의 분할입니다(4주차). 이사고게의 도구(류-종-종차-고유성-부수성)는 범주론의 모든 범주를 내부적으로 분석하는 데 사용됩니다.

여기서 사고 연습을 해봅시다. 관계 범주(4주차)에 포르피리오스식 나무를 적용할 수 있을까요? `관계'라는 최고류 아래에 `두 배/절반' 같은 수학적 관계와 `주인/노예' 같은 사회적 관계를 종차로 분할할 수 있을까요? 아리스토텔레스 자신은 이런 시도를 하지 않았습니다. 왜 실체 범주에서만 위계가 상세히 전개되고 다른 범주에서는 그렇지 않은지를 생각해 보면, 범주론이 실체 중심의 존재론임이 다시 한번 분명해집니다.

\subsection{우리가 만난 열린 문제들}

\subsubsection{범주의 상호 배타성 문제}

지난 주차들에서 반복적으로 마주친 문제가 있습니다 --- 하나의 항목이 두 범주에 동시에 속하는 것처럼 보이는 경우입니다.

지식(\gr{ἐπιστήμη})은 질 범주에 속합니다 --- 영혼의 상태(\gr{ἕξις})이니까요(4주차). 그런데 동시에 관계 범주에도 속합니다 --- 지식은 항상 무엇에 대한 지식이니까요(4주차). 물체(\gr{σῶμα})는 실체 범주의 종처럼 보이지만, 동시에 연속적 양의 사례이기도 합니다(3주차). 종차(`이성적인')는 본질적 술어이므로 실체에 가까워야 하지만, `어떤 종류의 것인지'를 말해주므로 질에 가깝습니다(7주차).

{\footnotesize
\begin{concepttable}{|l|c|c|l|l|}
\hline
\rowcolor{tableheadblue}
{\color{h1navy}\bfseries 항목} &
{\color{h1navy}\bfseries 범주 A} &
{\color{h1navy}\bfseries 범주 B} &
{\color{h1navy}\bfseries 해결 시도} &
{\color{h1navy}\bfseries 남은 문제} \\ \hline
지식 & 질 & 관계 & 류/종 구별 & 류가 관계적이면 종도? \\ \hline
물체 & 실체 & 양 & (명시적 해결 없음) & 이중 소속? \\ \hline
종차 & 실체? & 질? & 이사고게에서 별도 처리 & 사중 분류의 빈자리 \\ \hline
\end{concepttable}
}

\subsubsection{관계의 두 정의 --- 범주의 위기와 그 수습}

4주차에서 다루었던 관계의 두 정의 문제를 되돌아봅시다. 첫 번째 정의(6a36--37)는 너무 넓어서 머리, 손 같은 신체 부분 --- 즉 실체 --- 까지 관계에 포함시킬 위험이 있었습니다. 두 번째 정의(8a31--32)는 존재함이 곧 관계 맺고 있음인 것만을 관계로 인정하여 이 문제를 해결합니다. 이것은 범주론 내부에서 자기 수정이 일어나는 드문 사례입니다.\footnote{Matthew Duncombe, ``Aristotle's Two Accounts of Relatives in Categories 7,'' \gri{Phronesis} 60, no.\ 4 (2015): 436--461.}

\subsubsection{질 범주의 통일성 문제}

4주차에서 퍼스(Furth 1988)가 질 범주를 `기괴한 오합지졸'(monstrous motley horde)이라고 불렀다는 것을 보았습니다. 상태, 역량, 감각적 질, 모양 --- 이 네 종류가 왜 정확히 이 네 가지인지를 아리스토텔레스는 설명하지 않습니다. 스터트만(Studtmann 2003)과 아퀴나스의 체계적 도출 시도를 상기해 봅시다 --- 2천 년 뒤의 학자들이 여전히 아리스토텔레스의 분류를 정당화하려는 시도를 하고 있다는 것을 의미합니다.\footnote{Paul Studtmann, ``Aristotle's Category of Quality: A Regimented Interpretation,'' \gri{Apeiron} 36, no.\ 3 (2003): 205--227.}

\subsubsection{범주론과 형이상학의 긴장 --- `두 체계' 문제}

범주론에서 제1실체는 개별적 구체자(소크라테스, 이 말)입니다. 그러나 \citework{Met}{Z권}에서 제1실체는 형상(\gr{εἶδος}/\gr{μορφή}), 즉 사물을 그것이게 만드는 본질(\gr{τὸ τί ἦν εἶναι})입니다. 그레이엄(Graham 1987)은 이것을 `양립 불가능한 두 체계'로 해석했습니다.\footnote{Daniel W.\ Graham, \gri{Aristotle's Two Systems} (Oxford: Clarendon Press, 1987).} 웨딘(Wedin 2000)은 양립 가능하다고 봅니다 --- 범주론은 `어떤 존재자가 있는가'를 묻고, 형이상학은 `왜 그 존재자가 그런 종류의 것인가'를 묻는다는 것입니다.\footnote{M.\ V.\ Wedin, \gri{Aristotle's Theory of Substance: The Categories and Metaphysics Zeta} (Oxford: Oxford University Press, 2000).}

이 `두 체계' 문제는 우리의 강독 과정 전체에 배경으로 남아 있을 것입니다 --- 분석론 후서에서 정의와 논증의 관계를 다룰 때, 범주론의 실체 개념이 어떻게 변형되는지를 주시할 필요가 있습니다.

\subsection{고대 주석가들: 범주론은 무엇에 관한 책인가}

\subsubsection{암모니오스와 10개의 예비 물음}

암모니오스(\gr{Ἀμμώνιος}, c.\ 435--517/526)는 알렉산드리아 학파의 수장으로, 범주론을 읽기 위한 교육적 틀 자체를 설계했습니다.\footnote{암모니오스의 범주론 주석: CAG IV.4, ed.\ Busse. Eng.\ trans.: S.\ Marc Cohen and Gareth B.\ Matthews, \gri{Ammonius: On Aristotle's Categories} (Ithaca: Cornell University Press, 1991).} 그의 주석서는 학생들이 필기한 강의록(\gr{ἀπὸ τῆς φωνῆς}, `구술에서')의 형태로 전해집니다.

프로클로스(Proclus)가 원래 고안한 10개의 예비 물음(엘리아스 In Cat.\ 107.24--27의 증언)을 암모니오스가 체계화했습니다: (1) 철학 학파들의 이름은 어디서 오는가; (2) 아리스토텔레스의 저작 분류; (3) 어디서 공부를 시작하는가 --- 도구(\gr{ὄργανον})인 논리학에서; (4) 읽기 순서; (5) 철학의 유용성 --- `모든 것의 공통 원리로 올라가는 것'; (6) 학생의 자격; (7) 주석가의 자격 --- `아리스토텔레스에 대해 편파적이지 않아야 한다'; (8) 아리스토텔레스가 왜 의도적으로 난해하게 썼는가; (9) 문체적 특징; (10) 각 개별 저작에 대한 6--8개의 세부 물음.

이 구조는 러시아 인형 같습니다 --- 아리스토텔레스 일반에 대한 10개 물음, 그 다음 각 저작에 대한 6--8개 물음. 학생이 범주론의 첫 줄을 읽기 전에 수차례의 도입 강의를 거치는 셈입니다.

\subsubsection{세 가지 답변: 말인가, 개념인가, 사물인가}

범주론의 주제(\gr{σκοπός})를 둘러싼 고대의 논쟁은 범주론의 성격 자체를 결정하는 물음이었습니다. 알렉산드로스(fl.\ c.\ 200)의 존재론적 독해: 범주론은 사물(\gr{πράγματα})을 분류합니다. 포르피리오스의 의미론적 독해: 범주론은 의미하는 표현들(\gr{φωναὶ σημαντικαί})을 다룹니다 --- 표현 자체가 아니라, 표현이 사물을 의미하는 한에서. 이암블리코스(c.\ 245--325)의 지성적 독해(\gr{νοερὰ θεωρία}): 10범주는 감각 세계뿐 아니라 지성 세계에도 적용됩니다.

\subsubsection{심플리키오스의 대종합}

심플리키오스(c.\ 480--560)의 범주론 주석(CAG VIII)은 고대에 현존하는 가장 긴 주석서 가운데 하나입니다.\footnote{Eng.\ trans.: Michael Chase, \gri{Simplicius: On Aristotle's Categories 1--4}, Ancient Commentators on Aristotle (London: Duckworth, 2003).} 그의 해법은 세 입장의 통합이었습니다: `이 세 가지(말, 개념, 사물)는 원래 하나였으며, 나중에 분화되었다'(In Cat.\ 12.2). 이 삼중 통합은 현대 의미론의 `의미의 삼각형'(기호--사고--지시체)과 구조적으로 유사합니다.

\subsubsection{플로티노스의 다섯 가지 비판 --- 스승의 공격, 제자의 방어}

플로티노스의 엔네아데스 VI.1--3에서 제기된 다섯 가지 비판을 정리합시다.\footnote{Steven Strange, ``Porphyry and Plotinus' Metaphysics,'' in \gri{Aristotle Transformed}, ed.\ R.\ Sorabji (London: Duckworth, 1990), 955--974.} 첫째, 10범주는 감각 세계에만 적용됩니다. 둘째, 범주론의 `제1실체'는 진정한 실체가 아닙니다. 셋째, 질과 관계가 혼동됩니다. 넷째, 소유·위치·시간·장소는 불필요합니다. 다섯째, 언어가 양인 것은 소리인 한에서이지 의미적 본성에 있어서가 아닙니다.

포르피리오스의 대응은 정면 반박이 아니라 전략적 재해석이었습니다 --- 범주론을 형이상학이 아닌 의미론으로 재규정함으로써 플로티노스의 비판을 무력화한 것입니다. 이사고게의 유명한 첫머리 거부(보편자의 존재론적 지위에 답변 거부) 역시 같은 전략의 일환입니다.

여기서 사고 연습을 해봅시다. 포르피리오스의 전략은 성공적일까요? 범주론이 단순한 의미론이라면, 왜 아리스토텔레스는 제1실체에 존재론적 우선성(`다른 모든 것이 제1실체에 의존한다', 2b4--6)을 부여했을까요?

\subsubsection{덱시포스의 대화편}

덱시포스(\gr{Δέξιππος}, fl.\ c.\ 350)는 이암블리코스의 제자로, 범주론에 대한 유일한 대화편 형식의 고대 주석서를 남겼습니다.\footnote{Eng.\ trans.: John Dillon, \gri{Dexippus: On Aristotle Categories}, Ancient Commentators on Aristotle (London: Duckworth, 1990).} 심플리키오스는 덱시포스가 `선행자들에게 거의 아무것도 추가하지 않았다'(In Cat.\ 2.26--29)고 폄하했지만, 덱시포스의 저작은 포르피리오스와 이암블리코스의 소실된 주석서를 복원하는 데 대체 불가능한 자료입니다. 그의 접근법은 플로티노스의 비판과 아리스토텔레스의 입장이 생각만큼 멀지 않다는 것을 보여주려는 것이었습니다.

\subsection{우리가 획득한 것: 개념의 도구상자}

7주간의 강독을 통해 우리는 서양 철학의 기본 개념 어휘를 획득했습니다. 동음이의·동의·파생어(2주차), 주어에 대해 말해짐/주어 안에 있음(2주차), 제1실체/제2실체(2주차), 10범주(2주차), 반대자 수용(3주차), 이산적/연속적 양(3주차), 관계의 두 정의(4주차), 상태/상태됨(4주차), 네 가지 대립(5주차), 다섯 술어가능자(6주차), 보편자의 세 물음(6주차), 종차의 세 의미(7주차), 고유성의 네 의미(7주차), 분리 가능/불가능 부수성(8주차). 이 개념들은 명제론 이후의 오르가논 전체에서 계속 작동할 것입니다. 명제론에서는 범주론의 주어-술어 비대칭성이 명제의 논리적 구조로 발전하고, 분석론 전서에서는 술어화의 유형이 삼단논법의 타당성 조건을 결정합니다.

\subsection*{참고문헌}
\begin{references}

\refitem Dillon, John, trans. \gri{Dexippus: On Aristotle Categories}. Ancient Commentators on Aristotle. London: Duckworth, 1990.

\refitem Duncombe, Matthew. ``Aristotle's Two Accounts of Relatives in Categories 7.'' \gri{Phronesis} 60, no.\ 4 (2015): 436--461.

\refitem Evangeliou, Christos. \gri{Aristotle's Categories and Porphyry}. Leiden: Brill, 1996.

\refitem Graham, Daniel W. \gri{Aristotle's Two Systems}. Oxford: Clarendon Press, 1987.

\refitem Sorabji, Richard, ed. \gri{Aristotle Transformed: The Ancient Commentators and Their Influence}. London: Duckworth, 1990.

\refitem Strange, Steven. ``Porphyry and Plotinus' Metaphysics.'' In \gri{Aristotle Transformed}, edited by R.\ Sorabji, 955--974. London: Duckworth, 1990.

\refitem Studtmann, Paul. ``Aristotle's Category of Quality: A Regimented Interpretation.'' \gri{Apeiron} 36, no.\ 3 (2003): 205--227.

\refitem Wedin, M.\ V. \gri{Aristotle's Theory of Substance: The Categories and Metaphysics Zeta}. Oxford: Oxford University Press, 2000.

\end{references}

\end{document}
